% Hoeffding 1963: Theorem 3, its lemmas/comparisons, and Sections 5--6.
\chapter{Variance and Dependent-Sampling Concentration}
\label{sec:p15-variance-dependent-concentration}

The width of an interval and the variance of a variable contain different
information. An upper-tail bound can use a finite variance and an upper
bound without requiring any lower bound. Dependence requires an additional
argument: multiplying marginal exponential moments is justified by joint
independence, not by vanishing covariances.
This chapter reconstructs Hoeffding's variance theorem and its proof,
the variance comparisons, and the complete dependent-sampling constructions
in Sections 5 and 6 \cite[Thm.~3, Secs.~3--6]{Hoeffding1963BoundedSums}.
The common-range results used below are in
\cref{sec:p15-bounded-concentration-learning}.

\section{An upper bound and a finite variance}
\label{sec:p15-upper-variance}

\begin{hbdefinition}{def:p15-upper-variance-setting}{The centered upper-bounded setting}
Let \(n\geq1\) be an integer and let \(X_1,\ldots,X_n\) be jointly
independent real random variables with
\(EX_i=0\), \(v_i=EX_i^2<\infty\), and \(X_i\leq b\) almost surely,
where \(b>0\). Write
\[
 \overline X=n^{-1}\sum_{i=1}^nX_i,\qquad
 \sigma^2=n^{-1}\sum_{i=1}^nv_i.
\]
Thus \(\operatorname{Var}(\overline X)=\sigma^2/n\).
No common distribution or common individual variance is assumed.
For noncentered variables the condition is
\(X_i-EX_i\leq b\); centering must be performed separately for each variable.
\end{hbdefinition}

\begin{hblemma}{thm:p15-upper-moment-extremal}{proved=annex:p15.upper-moment-extremal}{The exact one-variable moment bound}
If \(b>0\), \(EX=0\), \(EX^2=v<\infty\), \(X\leq b\), and \(h>0\), then
\[
 Ee^{hX}\leq M_{b,v}(h):=
 \frac{b^2}{b^2+v}e^{-hv/b}+
 \frac{v}{b^2+v}e^{hb}.
\]
For \(v>0\), equality holds for the two-point distribution
\[
 \Pr\{X=-v/b\}=\frac{b^2}{b^2+v},\qquad
 \Pr\{X=b\}=\frac{v}{b^2+v}.
\]
For \(v=0\), \(X=0\) almost surely and \(M_{b,0}=1\).
This is Hoeffding's Lemma 2; the annex supplies a complete
Hermite-quadratic proof of the moment bound.
\end{hblemma}

\begin{hblemma}{thm:p15-variance-moment-concavity}{proved=annex:p15.variance-moment-concavity}{Averaging unequal variances}
For fixed \(c>0\), define
\[
 f_c(u)=\log\frac{e^{-cu}+ue^c}{1+u},\qquad u\geq0.
\]
Then \(f_c''(u)<0\) and \(f_c'(u)>0\), including the right derivatives
at zero. Also \(f_c(0)=0\) and \(f_c'(0)=e^c-1-c\).
Consequently, in \cref{def:p15-upper-variance-setting}, for every \(h>0\),
\[
 E e^{h\sum_iX_i}
 \leq M_{b,\sigma^2}(h)^n
 \leq \exp\!\left\{\frac{n\sigma^2}{b^2}(e^{hb}-1-hb)\right\}.
\]
The first comparison is the concavity argument of Hoeffding's Lemma 3;
the second is its tangent-line consequence.
\end{hblemma}

\begin{hbtheorem}{thm:p15-variance-tail-chain}{proved=annex:p15.variance-tail-chain}{Hoeffding's variance bound and its simpler consequences}
Assume \cref{def:p15-upper-variance-setting} and \(\sigma^2>0\).
For \(0<t<b\), put
\[
\begin{split}
 B_1(n,t;b,\sigma^2)
 &=
 \left\{
 (1+bt/\sigma^2)^{-(\sigma^2+bt)/(b^2+\sigma^2)}
 (1-t/b)^{-b(b-t)/(b^2+\sigma^2)}
 \right\}^{n},\\
 \lambda&=bt/\sigma^2,\qquad \tau=nt/b,\\
 h_1(\lambda)&=(1+\lambda^{-1})\log(1+\lambda)-1,\\
 h_2(\lambda)&=\frac{\lambda}{2(1+\lambda/3)},\qquad
 h_3(\lambda)=\frac12\operatorname{arsinh}(\lambda/2).
\end{split}
\]
Here \(\operatorname{arsinh}z=\log(z+\sqrt{1+z^2})\).
The source's equations (2.8)--(2.14) give
\[
 \Pr\{\overline X\geq t\}\leq B_1
 \leq B_2:=e^{-\tau h_1(\lambda)}
 \leq e^{-\tau h_2(\lambda)},\qquad
 B_2\leq e^{-\tau h_3(\lambda)}.
\]
In particular the Bernstein-form consequence is
\[
 \Pr\{\overline X\geq t\}
 \leq \exp\!\left\{-\frac{nt^2}{2(\sigma^2+bt/3)}\right\}.
\]
The Bennett-form bound \(B_2\) and its two simpler consequences remain
valid for every \(t>0\).

The name records Hoeffding's attribution. He credits the bound (2.9),
that is \(B_2\), to Bennett, as Bennett's inequality (8b), and notes that
Bennett's notation differs from his own. Bennett's first proof assumes
\(|X_i|\leq b\), with \(b\) playing the role of Bennett's \(M\); a second
proof, on pp.~42--43 of Bennett's paper, uses only \(X_i\leq b\)
\cite[p.~16]{Hoeffding1963BoundedSums}. That paper is item [2] of
Hoeffding's list, Bennett's 1962 article on probability inequalities for
sums of independent random variables in the same journal
\cite[p.~29]{Hoeffding1963BoundedSums}. Hoeffding also reports that
Bennett showed (2.12), the bound \(e^{-\tau h_1(\lambda)}\), to be better
than Bernstein's (2.13), the bound \(e^{-\tau h_2(\lambda)}\)
\cite[p.~17]{Hoeffding1963BoundedSums}; the comparison
\(h_1\geq h_2\) proved in the annex is that statement. The annex reaches
\(B_2\) from the tangent-line form of the moment bound. This is the
route Hoeffding offers as an alternative direct proof of Bennett's bound
\cite[p.~24]{Hoeffding1963BoundedSums}; in his own proof of Theorem~3,
(2.9) is instead read off from his equations (3.13)--(3.16).

At \(t=b\), the continuous extension of \(B_1\) is
\(\{\sigma^2/(b^2+\sigma^2)\}^n\); for \(t>b\) the probability is zero.
If \(\sigma^2=0\), every variable is zero almost surely and the probability
is zero for every \(t>0\). If the upper bound is \(b=0\), centering again
forces this degenerate case. An upper bound \(b<0\) is incompatible with
an integrable zero-mean variable.
For a two-sided deviation, apply the result also to \(-X_i\), which
requires its own upper bound, and add the two probabilities.
\end{hbtheorem}

\begin{hbcorollary}{thm:p15-variance-maximal}{proved=annex:p15.variance-maximal}{The source's independent partial-sum maximum}
Under \cref{def:p15-upper-variance-setting}, put
\(S_0=0\) and \(S_j=\sum_{i=1}^jX_i\). For \(h>0,t>0\),
\[
 \Pr\{\max_{0\leq j\leq n}S_j\geq nt\}
 \leq e^{-hnt}Ee^{hS_n}.
\]
All bounds and endpoint qualifications in
\cref{thm:p15-variance-tail-chain} therefore remain valid with
\(\{\overline X\geq t\}\) replaced by this maximum event.
The variance parameter remains \(n^{-1}\sum_i EX_i^2\).
The independence premise is retained; this corollary does not replace
it by a statement about the unconditional variances of dependent
martingale increments.

Hoeffding states this strengthening, his (2.17), as a consequence of a
theorem of Doob, and calls the stronger result analogous to an
inequality of Kolmogorov, for which he refers to Feller
\cite[pp.~17--18]{Hoeffding1963BoundedSums}. His bracketed reference
for Doob reads [7, p.~314], although item [7] of his reference list is
Hardy, Littlewood and P\'olya's \emph{Inequalities} and Doob's
\emph{Stochastic Processes} (1953) is item [4]
\cite[p.~29]{Hoeffding1963BoundedSums}. The annex proves the corollary
directly by a first-crossing decomposition and does not use Doob's
theorem. The analogy is that both results bound the probability that a
running sum ever reaches a level by a moment of the final sum alone. The
same decomposition with \(e^{hS_n}\) replaced by \(S_n^2\) gives
Kolmogorov's inequality: for jointly independent centered \(X_i\) with
finite variances, with no upper bound assumed, and for \(x>0\),
\[
 \Pr\{\max_{1\leq j\leq n}|S_j|\geq x\}\leq\frac{ES_n^2}{x^2}.
\]
\end{hbcorollary}

\begin{hbtheorem}{thm:p15-variance-bound-comparison}{proved=annex:p15.variance-bound-comparison}{Exact comparison and the role of known variance}
For \(0<t<b\) and \(\sigma^2>0\), define
\[
 v=\frac{bt}{bt+\sigma^2},\quad w=\frac tb,\quad
 \rho(x)=\frac{(1-x)\log(1-x)+x-x^2/2}{x^2}\quad(0<x<1).
\]
Then
\[
 \frac{B_1}{B_2}=e^{-n\phi(v,w)},\qquad
 \phi(v,w)=\frac{w\{\rho(v)+\rho(w)\}}{v^{-1}+w^{-1}-1},
 \qquad
 \rho(x)=\sum_{j=1}^{\infty}\frac{x^j}{(j+1)(j+2)}.
\]
In particular
\[
 \frac{w(v+w)}{6(v^{-1}+w^{-1}-1)}
 <\phi(v,w)<
 \frac{w(v+w)}{2(v^{-1}+w^{-1}-1)}
 <\frac{t^3}{2b\sigma^2}.
\]
Thus \(e^{-nt^3/(2b\sigma^2)}<B_1/B_2<1\); the left inequality is the
source's (3.15). Section~3 of the source measures such comparisons by a
ratio: two bounds count as appreciably different when their ratio is
not close to one \cite[p.~18]{Hoeffding1963BoundedSums}. Read with that
criterion, (3.15) says that \(B_1\) is not appreciably smaller than
\(B_2\) whenever \(nt^3/(2b\sigma^2)\) is small.
For fixed positive \(b,\sigma^2\), as \(t\downarrow0\),
\(\phi(v,w)=t^3/(6b\sigma^2)+O(t^4)\).

Hoeffding states that when \(bt/\sigma^2\) is small, so that \(v\) is
close to \(bt/\sigma^2\), the ratio \(B_1/B_2\) is approximately
\(\exp\{-nt^3/(6b\sigma^2)\}\) \cite[p.~20]{Hoeffding1963BoundedSums}.
The precise form of this regime is as follows. With \(\lambda=bt/\sigma^2\)
as above, put \(M=\max\{\lambda,t/b\}\). If \(M<1\), then
\[
 \frac1{1+\lambda}
 <\frac{\phi(v,w)}{t^3/(6b\sigma^2)}
 <\frac{2-M}{2(1-M)^2}.
\]
So when both \(bt/\sigma^2\) and \(t/b\) are small, the exponent
\(n\phi(v,w)\) is the product of \(nt^3/(6b\sigma^2)\) and a factor
\(1+O(M)\), and by the ratio criterion \(B_1\) is appreciably smaller
than \(B_2\) only when \(nt^3/(6b\sigma^2)\) is not small. Since
\(t/b=\lambda\sigma^2/b^2\), smallness of \(\lambda\) alone suffices when
\(\sigma^2\leq b^2\); this holds, for instance, under a two-sided
range \(a\leq X_i\leq b\) with \(-a\leq b\), because then
\(\sigma^2\leq-ab\leq b^2\). When \(\sigma^2>b^2\) it does not suffice.
With \(w=t/b\) held fixed and \(\lambda\downarrow0\),
\[
 \frac{\phi(v,w)}{t^3/(6b\sigma^2)}\longrightarrow\frac{6\rho(w)}w
 >1+\frac w2,
\]
so the exponent is asymptotically \((nt^2/\sigma^2)\rho(t/b)\) rather
than \(nt^3/(6b\sigma^2)\). At \(w=1/2\) the limiting factor is
\(18-24\log2\approx1.364\).

Both \(h_1(\lambda)\) and \(h_2(\lambda)\) equal
\(\lambda/2-\lambda^2/6+O(\lambda^3)\) as \(\lambda\downarrow0\).
Explicitly, for \(0<\lambda\leq1\),
\[
 h_1(\lambda)=\sum_{k=1}^\infty\frac{(-1)^{k+1}\lambda^k}{k(k+1)},
 \qquad
 0<h_1(\lambda)-\frac\lambda2+\frac{\lambda^2}6<\frac{\lambda^3}{12},
\]
and for every \(\lambda>0\),
\(h_2(\lambda)-\lambda/2+\lambda^2/6=\lambda^3/\{18(1+\lambda/3)\}\).

The bound \(B_1\) is strictly increasing in \(\sigma^2\), as Hoeffding
states for (2.8) \cite[p.~20]{Hoeffding1963BoundedSums}: for
\(0<\sigma_1^2<\sigma_2^2\), the value at \(\sigma_1^2\) is strictly
smaller than the value at \(\sigma_2^2\).
If also \(a\leq X_i\leq b\), where \(a<0<b\), then
\(\sigma^2\leq-ab\). Replacing \(\sigma^2\) by \(-ab\) in \(B_1\)
gives exactly the common-range exponential bound from
\cref{thm:p15-conc-common-range}, with normalized mean
\(-a/(b-a)\) and normalized displacement \(t/(b-a)\).
This bridge assumes the separately centered means of the present theorem;
it does not add a common-mean requirement to the source's Theorem 1.

The Bennett-form bound \(B_2\) is also nondecreasing in \(\sigma^2\).
Under the same premises \(a\leq X_i\leq b\), \(a<0<b\), the same
substitution in \(B_2\), the source's (2.9), gives for every \(t>0\)
\[
 \Pr\{\overline X\geq t\}\leq
 \exp\!\left\{-\frac{nt}{b}
 \left[\Bigl(1-\frac at\Bigr)\log\Bigl(1-\frac ta\Bigr)-1\right]\right\},
\]
a bound that involves \(a\) and \(b\) but no variance. For \(0<t<b\) it is
never smaller than the common-range bound obtained above from \(B_1\):
their ratio is \(e^{-n\phi(v,w)}\) with \(v=t/(t-a)\) and \(w=t/b\).
Hoeffding adds that this variance-free form is no simpler than the
smaller bound that corresponds to his (2.1)
\cite[p.~20]{Hoeffding1963BoundedSums}; that remark is an assessment of
form, and the mathematical content is the displayed order.
In the first calculation of \cref{thm:p15-dependent-concentration-examples},
\(a=-1/4\), \(b=1\) and \(-ab=1/4=\sigma^2\), so neither substitution
changes anything: the two variance-free bounds are \(2^{16}/3^{12}\) and
\(e^5/2^{10}\), and the first is the common-range bound with \(p=1/5\),
\(q=2/5\).
\end{hbtheorem}

\begin{hbtheorem}{thm:p15-variance-optimality-boundary}{proved=annex:p15.variance-optimality-boundary}{What exponential optimality does and does not mean}
For fixed integer \(n\geq1\), \(b,\sigma^2>0\), and \(0<t<b\),
\(B_1\) is the minimum exponential-Markov
bound for the independent equal-variance two-point law in
\cref{thm:p15-upper-moment-extremal}. That law also attains the universal
moment bound at every positive \(h\). At \(t=b\) its tail probability
equals \(B_1\).

This does not make \(B_1\) the least possible uniform bound on the actual
tail for \(0<t<b\). For \(n=1,b=1,\sigma^2=1,t=1/2\),
\[
 B_1=\frac{2}{3^{3/4}}>\frac45.
\]
The sharp uniform tail bound at those values is \(4/5\), attained by
the law with masses \(4/5,1/5\) at \(1/2,-2\).

More generally, let \(W\) be real, let \(J\geq1\) be finite, and choose
real \(c_j,\alpha_j\) with \(Ee^{\alpha_jW}<\infty\).
If \(Q(s)=\sum_{j=1}^Jc_je^{\alpha_js}\) satisfies
\(Q(W)\geq0\) almost surely and \(Q(W)\geq1\) on \(\{W\geq a\}\),
then
\[
 \Pr\{W\geq a\}\leq\sum_{j=1}^Jc_jEe^{\alpha_jW}.
\]
This is the source footnote's more general exponential-majorant route.
A negative coefficient requires a lower moment bound, or the exact
moment, when bounding its contribution from above.

Among reductions of the form \(e^{-\tau H(\lambda)}\) obtained uniformly
from \(B_1\) for all admissible positive \(b,\sigma^2,t\), the largest
possible exponent function is \(H=h_1\). This last assertion concerns
that reduction of an exponential bound, not exact tail optimization.
\end{hbtheorem}

\begin{hbtheorem}{thm:p15-variance-upper-bound-crucial}{proved=annex:p15.variance-upper-bound-crucial}{Without the upper bound, the one-sided Chebyshev bound is sharp}
Fix an integer \(n\geq1\), \(t>0\) and \(\sigma^2>0\). Let
\(\mathcal C(n,\sigma^2)\) be the class of jointly independent real
\(X_1,\ldots,X_n\) with \(EX_i=0\), \(EX_i^2<\infty\) and
\(n^{-1}\sum_iEX_i^2=\sigma^2\): the premises of
\cref{def:p15-upper-variance-setting} with the condition \(X_i\leq b\)
deleted. Then
\[
 \sup_{\mathcal C(n,\sigma^2)}\Pr\{\overline X\geq t\}
 =\frac{\sigma^2}{\sigma^2+nt^2}
 =\frac1{1+nt^2/\sigma^2},
\]
the right side of the one-sided Chebyshev bound (1.6) with \(\mu=0\).
The supremum is attained by \(X_2=\cdots=X_n=0\) and
\[
 \Pr\{X_1=nt\}=\frac{\sigma^2}{\sigma^2+nt^2},\qquad
 \Pr\{X_1=-\sigma^2/t\}=\frac{nt^2}{\sigma^2+nt^2}.
\]
Consequently no upper bound for \(\Pr\{\overline X\geq t\}\) that holds
on \(\mathcal C(n,\sigma^2)\) and depends only on \(n\), \(t\) and
\(\sigma^2\) is smaller than (1.6). This is the second half of
Hoeffding's footnote to Theorem~1, which says that in Theorem~3 the
assumption \(X_i\leq b\) is crucial to getting any improvement over
(1.6) \cite[p.~15 n.~4]{Hoeffding1963BoundedSums}; the bound (1.6) is
displayed at \cite[p.~13]{Hoeffding1963BoundedSums}.

The attaining law satisfies \(X_i\leq nt\). Hence, for every \(b\geq nt\),
the supremum of \(\Pr\{\overline X\geq t\}\) over the variables of
\cref{def:p15-upper-variance-setting} is still \(\sigma^2/(\sigma^2+nt^2)\);
a bound that improves on (1.6) for all such variables can exist only
when \(b<nt\). For fixed \(b\) and \(t\) this holds for all large \(n\),
and there \(B_1\) and \(B_2\) of \cref{thm:p15-variance-tail-chain}
decrease exponentially in \(n\), whereas (1.6) decreases only like
\(\sigma^2/(nt^2)\). The case \(n=1\), \(\sigma^2=1\), \(t=1/2\) is the
law with masses \(4/5,1/5\) at \(1/2,-2\) in
\cref{thm:p15-variance-optimality-boundary}.
The companion half of the same footnote, for Theorem~1 without the
upper range, is the Markov attainment in
\cref{thm:p15-conc-benchmarks}.
\end{hbtheorem}

\section{Scope, names and priority in Hoeffding's account}
\label{sec:p15-w02c-p15-017-hoeffding-scope}

This section collects what Hoeffding says about the reach of his bounds
and about their history, and proves the mathematical claims among those
remarks. The notation of \cref{sec:p15-bounded-concentration-learning}
is used for the common-range bounds: \(A_1,A_2,A_3\) as in
\cref{thm:p15-conc-common-comparison}, and \(p\) for the average mean
that the source writes \(\mu\).

\emph{The comparison bounds and their names.} Among the known bounds
that Hoeffding lists for comparison are two moment bounds. He calls the two-sided bound (1.5),
\(\Pr\{|\overline X-\mu|\geq t\}\leq\sigma^2/(nt^2)\), the
Bienaym\'e--Chebyshev inequality, and the one-sided bound (1.6),
\(\Pr\{\overline X-\mu\geq t\}\leq1/(1+nt^2/\sigma^2)\), Chebyshev's
inequality \cite[p.~13]{Hoeffding1963BoundedSums}. A footnote to (1.6)
records that this one-sided inequality has been credited to several
authors, and that Chebyshev appears to have been the first to announce
an inequality implying it, as one instance of a general class of
inequalities \cite[p.~13 n.~2]{Hoeffding1963BoundedSums}. The
reference is item [14] of Hoeffding's list, Tchebichef's note ``Sur les
valeurs limites des int\'egrales'' in the \emph{Journal de
Math\'ematiques Pures et Appliqu\'ees}, ser.~2, vol.~19 (1874),
pp.~157--160 \cite[p.~30]{Hoeffding1963BoundedSums}. Elsewhere in this
part the pair is called the Chebyshev bounds
(\cref{thm:p15-conc-benchmarks,thm:p15-w02b-p15-063-chebyshev-comparison});
the two names above are the source's own, and the attribution history is
the one its footnote reports.

\emph{Surveys.} Hoeffding names three surveys of inequalities of this
type: Godwin [6], Savage [13] and Bennett [2]. He adds that Bennett also
obtained new inequalities, (2.12) among them, and compared
several bounds with one another in an instructive way \cite[p.~14]{Hoeffding1963BoundedSums}.
In his reference list these are Godwin, ``On generalizations of
Tchebychef's inequality'', \emph{Journal of the American Statistical
Association} 50 (1955), pp.~923--945; Savage, ``Probability inequalities
of the Tchebycheff type'', \emph{Journal of Research of the National
Bureau of Standards} 65B (1961), pp.~211--222; and Bennett's 1962 paper
in the \emph{Journal of the American Statistical Association} 57,
pp.~33--45 \cite[pp.~29--30]{Hoeffding1963BoundedSums}. The bound (2.12)
is \(e^{-\tau h_1(\lambda)}\) of \cref{thm:p15-variance-tail-chain}. The
surveys are cited here as Hoeffding reports them; nothing in this chapter
is drawn from them.

\emph{The exponential method.} Every bound of Hoeffding's Theorems~1--3
comes from one step: the indicator of \(\{S-ES\geq nt\}\) is at most
\(e^{h(S-ES-nt)}\) for \(h>0\), which gives his (1.7), and for independent
summands the exponential moment factorizes as in his (1.8)
\cite[p.~14]{Hoeffding1963BoundedSums}. In this handbook the step is the
last display of \cref{thm:p15-conc-exponential}. Hoeffding does not
claim the method. He says that it has often been used, and apparently
first by S.~N.~Bernstein \cite[p.~14]{Hoeffding1963BoundedSums}; he gives
no reference to Bernstein at that point. Where this part speaks of
Hoeffding's method, as in \cref{sec:p15-w02b-p15-063-fixed-sample}, it
means the method as Hoeffding applies it, not its origin.

\emph{Priority for the binomial case.} For the binomial case (2.5),
independent summands each taking the values \(0\) and \(1\) with common
mean \(p\), Hoeffding reports that the inequalities of his Theorem~1 had
already been derived by Okamoto [11], with one exception: the second
bound (2.2) when \(p<1/2\). He reports further that the binomial case of (2.1)
is implicit in Chernoff's paper (his reference [3], Theorem~1 and
Example~5) \cite[pp.~15--16]{Hoeffding1963BoundedSums}. His list
gives Okamoto, ``Some inequalities relating to the partial sum of binomial
probabilities'', \emph{Annals of the Institute of Statistical
Mathematics} 10 (1958), pp.~29--35, and Chernoff, ``A measure of
asymptotic efficiency for tests of a hypothesis based on the sum of
observations'', \emph{Annals of Mathematical Statistics} 23 (1952),
pp.~493--507 \cite[pp.~29--30]{Hoeffding1963BoundedSums}. In the notation
of \cref{thm:p15-conc-common-range}, (2.1) is the bound
\(e^{-nD(p+t\Vert p)}\) and (2.2) with \(p<1/2\) is
\(e^{-nt^2g(p)}\) with the logarithmic branch of \(g\). The theorem itself
assumes only independence and \(0\leq X_i\leq1\); in the binomial case the
minimum over \(h\) of the exponential bound equals (2.1), by part~2 of
\cref{thm:p15-w02b-p15-063-exponential-optimality}.

\emph{Exact bounds, not approximations.} The footnote that defines the
best possible bound, used in \cref{thm:p15-conc-common-range}, goes on
to say that several authors, Bahadur and Rao [1] in particular, had
studied approximations to that probability built from the right side of
(1.7) after minimization over \(h\), and that their paper lists earlier
work. Hoeffding then sets his own aim apart: he seeks exact bounds, and
does not approximate the probability \cite[p.~14 n.~3]{Hoeffding1963BoundedSums}. Item [1] is
Bahadur and Rao, ``On deviations of the sample mean'', \emph{Annals of
Mathematical Statistics} 31 (1960), pp.~1015--1027
\cite[p.~29]{Hoeffding1963BoundedSums}. Accordingly every inequality of
Theorems~1--3, and every bound of
\cref{sec:p15-bounded-concentration-learning} and of this chapter, holds
at each finite \(n\) under its stated hypotheses. The normal limit in
\cref{thm:p15-conc-benchmarks} is of the other kind: it describes the
probability as \(n\to\infty\) and bounds it at no finite \(n\).

\begin{hbtheorem}{thm:p15-w02c-p15-017-nonpositive-deviation}{proved=annex:p15-w02c-p15-017-nonpositive-deviation}{No nontrivial bound for a nonpositive deviation}
Let \(n\geq1\) and \(t\leq0\).
\begin{enumerate}
\item For every \(p\in[0,1]\), the constant variables \(X_i=p\) are
independent, satisfy \(0\leq X_i\leq1\), have average mean \(p\), and give
\(\Pr\{\overline X-p\geq t\}=1\). The same holds for constant variables
\(X_i=c_i\in[a_i,b_i]\) under the ranges of
\cref{thm:p15-conc-unequal-ranges}, with \(S_n=0\).
\item Fix \(b>0\) and \(\sigma^2>0\). For each \(q\) with
\(0<q\leq b^2/(b^2+n\sigma^2)\), let \(X_2=\cdots=X_n=0\) and let \(X_1\)
take the value \(\sigma\sqrt{nq/(1-q)}\) with probability \(1-q\) and the
value \(-\sigma\sqrt{n(1-q)/q}\) with probability \(q\). These variables
satisfy \cref{def:p15-upper-variance-setting} with this \(b\) and average
variance \(\sigma^2\), and \(\Pr\{\overline X\geq t\}\geq1-q\). Hence the
supremum of \(\Pr\{\overline X\geq t\}\) over that setting is one.
\end{enumerate}
So for \(t\leq0\) no upper bound smaller than one depends only on \(t\),
\(n\), the ranges and the mean, or on \(t\), \(n\), \(b\) and
\(\sigma^2\). This is why Hoeffding assumes \(t>0\): for \(t\leq0\) no
nontrivial upper bound exists under his assumptions
\cite[p.~13]{Hoeffding1963BoundedSums}.
\end{hbtheorem}

At \(q=b^2/(b^2+n\sigma^2)\) the law of \(X_1\) in part~2 is the
two-point law of \cref{thm:p15-upper-moment-extremal} for the variance
\(n\sigma^2\), with upper value \(b\). In part~1 the supremum is
attained; in part~2 it is approached as \(q\downarrow0\).

\begin{hbtheorem}{thm:p15-w02c-p15-017-markov-comparison}{proved=annex:p15-w02c-p15-017-markov-comparison}{When the common-range bounds improve on Markov's bound}
Assume the hypotheses of \cref{thm:p15-conc-common-range}: independent
\(0\leq X_i\leq1\), average mean \(p\in(0,1)\) and \(0<t<1-p\). Since
\(\overline X\geq0\), Markov's inequality gives
\[
 \Pr\{\overline X-p\geq t\}=\Pr\{\overline X\geq p+t\}\leq\frac p{p+t}.
\]
\begin{enumerate}
\item If \(n=1\), then \(p/(p+t)<A_1\leq A_2\leq A_3\).
\item For every \(n\geq1\), if \(nt(p+t)\leq p(1-p)\), then
\(p/(p+t)<A_1\leq A_2\leq A_3\).
\item If \(2npt>1\), then \(A_1\leq A_2\leq A_3<p/(p+t)\).
\item At the endpoint \(t=1-p\), where \(A_1\) extends to \(p^n\), the two
bounds agree for \(n=1\) and \(p^n<p\) for \(n\geq2\).
\end{enumerate}
Thus, for fixed \(p\) and \(n\), the three bounds are larger than Markov's
for all deviations \(t\) with \(nt(p+t)\leq p(1-p)\); for fixed \(p\) and
\(t\), all three are smaller once \(n>1/(2pt)\); and for a single
summand none of them improves on Markov's bound at any \(t\) with
\(0<t<1-p\).
\end{hbtheorem}

Hoeffding's footnote to Theorem~1 opens by saying that its remarks are
based on a referee's comments. Its first remark is the attainment of
\(\mu/(\mu+t)\) when the upper bound \(X_i\leq1\) is dropped
(\cref{thm:p15-conc-benchmarks}), from which it concludes that
boundedness on both sides is crucial to any improvement over Markov's
bound; it then adds, in parentheses, that the gain appears only once
neither \(n\) nor \(t\) is too small \cite[p.~15 n.~4]{Hoeffding1963BoundedSums}.
\Cref{thm:p15-w02c-p15-017-markov-comparison} states what ``not too
small'' amounts to for the bounds of Theorem~1: part~1 shows that \(n=1\)
is always too small, part~2 that small \(t\) is too small at every \(n\),
and part~3 gives a sufficient size of \(n\) for each fixed \(t\). The
same footnote's remark about Theorem~3 is
\cref{thm:p15-variance-upper-bound-crucial}. Part~1 agrees with part~4 of
\cref{thm:p15-w02b-p15-063-exponential-optimality}, where \(p/(p+t)\) is
shown to be the least upper bound for one summand.

\begin{hbexample}{def:p15-w02c-p15-017-markov-example}{A fair-coin threshold}
Take \(p=1/2\) and \(t=1/4\), so Markov's bound is \(2/3\). Here
\(A_2=A_3=e^{-n/8}\) and
\[
 A_1=e^{-nD(3/4\Vert1/2)}=\Bigl(\frac2{3^{3/4}}\Bigr)^n .
\]
Part~2 of the theorem applies only at \(n=1\), where
\(A_1=2/3^{3/4}\approx0.877\). Part~3 applies from \(n=5\). The exact
crossing is earlier: \(A_1>2/3\) for \(n\leq3\) (at \(n=3\),
\(A_1=8/3^{9/4}\approx0.675\)), \(A_3>2/3\) for \(n\leq3\)
(at \(n=3\), \(e^{-3/8}\approx0.687\)), and at \(n=4\)
\(A_1=16/27\approx0.593\) and \(A_3=e^{-1/2}\approx0.607\) are both below
\(2/3\). So in this case all three bounds improve on Markov's exactly
when \(n\geq4\).
\end{hbexample}

\begin{hbcorollary}{thm:p15-w02c-p15-017-section-three-reading}{proved=annex:p15-w02c-p15-017-section-three-reading}{Reading the ratio estimates for the first two common-range bounds}
Under the hypotheses of \cref{thm:p15-conc-common-comparison} with
\(p\geq1/2\), let \(L\) and \(U\) be the left and right members of the
two-sided estimate of \(A_1/A_2\) there: for \(p>1/2\),
\[
 L=\exp\!\left\{-\frac{nt^3}{3(1-p)(1-p-t)}\right\},\qquad
 U=\exp\!\left\{-\frac{(2p-1)nt^3}{6p^2(1-p)^2}\right\},
\]
and for \(p=1/2\), \(L=\exp\{-4nt^4/(3(1-4t^2))\}\) and
\(U=e^{-4nt^4/3}\). Then \(0<L<A_1/A_2<U<1\), and
\[
 A_1<U\,A_2,\qquad 0<\frac{A_2-A_1}{A_2}<1-L .
\]
So if \(U\) is not close to one, the first bound is smaller than the
second by at least the factor \(U\); and if \(L\) is close to one, the
second bound exceeds the first by less than the fraction \(1-L\) of
itself.
\end{hbcorollary}

Hoeffding opens his Section~3 by observing that in Theorem~1 each of the
three bounds trades strength for simplicity against its predecessor, and
that the second bound of Theorem~3 does the same against the first. His
question is when a simpler bound comes near a more complicated one and
when the complicated one is appreciably better, and he counts two bounds
as appreciably different when their ratio is not close to one
\cite[p.~18]{Hoeffding1963BoundedSums}. After (3.6) and (3.7) he draws
the conclusion stated in the corollary: right-hand sides far from one
mean that the first bound is appreciably smaller, and left-hand sides near
one mean that little is lost by using the simpler second bound instead
\cite[p.~19]{Hoeffding1963BoundedSums}.
The corollary is that reading in the form of inequalities. The same
criterion applied to \(B_1/B_2\) is in
\cref{thm:p15-variance-bound-comparison}. For \(p<1/2\) the analogous
statement uses the lower estimate \(e^{-n(1-2p-t)^2/p}\leq A_1/A_2\leq1\)
of \cref{thm:p15-conc-common-comparison}: when that lower estimate is
close to one, so is the ratio. Hoeffding's p.~19 conclusion is stated for
(3.6) and (3.7) only.

The two members need not settle the question between them. If \(L\) is
not close to one while \(U\) is, the estimate leaves open whether the
first bound is appreciably better. For \(p>1/2\) the two exponents differ
by the factor
\[
 \frac{2p^2(1-p)}{(2p-1)(1-p-t)},
\]
which exceeds one, and is large when \(p\) is near \(1/2\).

\begin{hbexample}{def:p15-w02c-p15-017-ratio-example}{Both readings at one mean and deviation}
Take \(p=3/4\) and \(t=1/8\), so \(t<1-p=1/4\). Then
\(U=e^{-n/216}\) and \(L=e^{-n/48}\). At \(n=4\), \(L=e^{-1/12}\), so
\(A_2\) exceeds \(A_1\) by less than one twelfth of \(A_2\): the simpler
bound is almost as good. At \(n=432\), \(U=e^{-2}\), so
\(A_1<e^{-2}A_2<0.14\,A_2\): the first bound is appreciably better.
\end{hbexample}

\section{Convex averages of independent blocks}
\label{sec:p15-dependent-blocks}

\begin{hblemma}{thm:p15-concentration-convex-mixture}{proved=annex:p15.concentration-convex-mixture}{A pointwise convex average transfers moment bounds}
Let \(T=\sum_{j=1}^Np_jT_j\), with deterministic \(p_j\geq0\),
\(\sum_jp_j=1\), and integrable real \(T_j\) on one probability space.
The \(T_j\) may be dependent. For every \(h>0\),
\[
 E e^{h(T-ET)}
 \leq\sum_jp_j E e^{h(T_j-ET_j)},\qquad
 \Pr\{T-ET\geq t\}
 \leq e^{-ht}\sum_jp_j E e^{h(T_j-ET_j)}.
\]
Terms with zero weight are omitted. The comparisons allow infinite right
sides. A common finite moment
bound for all the centered \(T_j\) therefore transfers to \(T\).
This is an average of variables on the same sample, not an assumption
that a random mixture selects one independent component.

Hoeffding's own form, his (5.2), is uncentered: for every real \(t\) and
\(h>0\),
\[
 \Pr\{T\geq t\}\leq\sum_{j=1}^Np_jEe^{h(T_j-t)}
\]
\cite[p.~24]{Hoeffding1963BoundedSums}. The centered display above is
(5.2) applied to the variables \(T_j-ET_j\), whose weighted average is
\(T-ET\). The two forms agree when every \(ET_j\) equals \(ET\). When
the means differ, the centered form uses the same deviation \(t\) for
every component, while in (5.2) the effective deviation of \(T_j\) above
its own mean is \(t-ET_j\), different for each \(j\).

If the \(T_j\) are identically distributed, or if a bound for
\(Ee^{h(T_j-t)}\) does not depend on \(j\), the bound obtained for
\(\Pr\{T\geq t\}\) is also a bound for each \(\Pr\{T_j\geq t\}\).
Hoeffding remarks for this reason that the bounds produced by the
method will be rather crude, though they may be useful
\cite[p.~24]{Hoeffding1963BoundedSums}. In general the right side of
(5.2) is at least \(\min_jEe^{h(T_j-t)}\): an average of the components
is bounded no better than its best single component, even when the
average is much more concentrated than any component.
\end{hblemma}

\begin{hbdefinition}{def:p15-ordered-u-statistic}{Ordered tuples and block averages}
For integers \(n\geq r\geq1\), put
\((n)_r=n(n-1)\cdots(n-r+1)\), and set \((m)_0=1\) for every integer
\(m\geq0\). Let \(X_1,\ldots,X_n\) be independent
random elements in a common measurable space \(\mathcal X\), and let
\(g:\mathcal X^r\to\mathbb R\) be a measurable function of \(r\)
arguments, integrable for every distinct-index tuple. Define
\[
 U=\frac1{(n)_r}\sum_{i_1,\ldots,i_r\ {\rm distinct}}
                   g(X_{i_1},\ldots,X_{i_r}),\qquad k=\lfloor n/r\rfloor.
\]
No symmetry of \(g\) is required: the average uses ordered tuples.
For a permutation \(\pi\) of \(\{1,\ldots,n\}\), put
\[
 V_\pi=\frac1k\sum_{j=1}^k
 g(X_{\pi((j-1)r+1)},\ldots,X_{\pi(jr)}).
\]
Each \(V_\pi\) averages \(k\) mutually independent block variables.
\end{hbdefinition}

\begin{hbtheorem}{thm:p15-one-sample-u-concentration}{proved=annex:p15.one-sample-u-concentration}{One-sample permutation representation}
In \cref{def:p15-ordered-u-statistic},
\[
 U=\frac1{n!}\sum_\pi V_\pi.
\]
If \(a\leq g\leq b\), \(a<b\), then for \(t>0\)
\[
 \Pr\{U-EU\geq t\}\leq e^{-2kt^2/(b-a)^2}.
\]
If the \(X_i\) are identically distributed, set
\(\mu=Eg(X_1,\ldots,X_r)=EU\).
For \(0\leq g\leq1\), all three common-range bounds and their endpoint
qualifications in \cref{thm:p15-conc-common-range} transfer with \(n=k\).
If \(g-\mu\leq b_0\), \(b_0>0\), and
\(\sigma_g^2=\operatorname{Var}(g(X_1,\ldots,X_r))<\infty\), the bounds
in \cref{thm:p15-variance-tail-chain} transfer with
\(n=k,b=b_0,\sigma^2=\sigma_g^2\).
The variance here is that of one kernel value, not of \(U\).
\end{hbtheorem}

\begin{hbtheorem}{thm:p15-two-sample-u-concentration}{proved=annex:p15.two-sample-u-concentration}{Two-sample \(U\) statistics and their permutation representation}
Let all \(X_1,\ldots,X_m,Y_1,\ldots,Y_n\) be mutually independent,
\(m\geq r\geq1,n\geq s\geq1\), and let \(g\) be a measurable integrable
kernel of \(r\) first-sample and \(s\) second-sample arguments.
All \(X_i\) take values in one measurable space and all \(Y_j\) in
another; the kernel is defined on the corresponding product. Define
\[
 U=\frac1{(m)_r(n)_s}
 \sum_{\substack{i_1,\ldots,i_r\ {\rm distinct}\\
                 j_1,\ldots,j_s\ {\rm distinct}}}
 g(X_{i_1},\ldots,X_{i_r},Y_{j_1},\ldots,Y_{j_s}),\quad
 k=\min(\lfloor m/r\rfloor,\lfloor n/s\rfloor).
\]
A random variable of this form, the source's (5.8), is a
\emph{two-sample \(U\) statistic} \cite[pp.~25--26]{Hoeffding1963BoundedSums}.
\index{U statistic!two-sample}
For permutations \(\pi,\eta\), let \(V_{\pi,\eta}\) average the \(k\)
kernel values formed from the successive disjoint \(r\)-blocks in
\(\pi\) and successive disjoint \(s\)-blocks in \(\eta\). Then
\[
 U=\frac1{m!n!}\sum_{\pi,\eta}V_{\pi,\eta}.
\]
If \(a\leq g\leq b\), \(a<b\), then for every \(t>0\) the upper-tail
bound for \(U-EU\) is \(e^{-2kt^2/(b-a)^2}\).
If each sample is identically distributed within itself, the common-range
and variance-sensitive conclusions of the preceding theorem hold with
this \(k\) and the mean and variance of one mixed kernel.
The two sample distributions need not be equal.
\end{hbtheorem}

\begin{hbtheorem}{thm:p15-repeated-index-concentration}{proved=annex:p15.repeated-index-concentration}{Repeated indices reduce to an explicit kernel average}
Let \(n\geq r\geq1\), with independent \(X_i\) in a common measurable
space and a real measurable kernel \(g\) on its \(r\)-fold product.
Define
\[
 W=\frac1{n^r}\sum_{i_1,\ldots,i_r=1}^n
                       g(X_{i_1},\ldots,X_{i_r}).
\]
Assume all displayed terms integrable, including repeated arguments.
For a partition \(\mathcal P\) of \(\{1,\ldots,r\}\) into \(j\) nonempty
blocks, fix a labeling of its blocks and let \(\beta_{\mathcal P}(\ell)\)
be the label of the block containing \(\ell\). Set
\[
 g^*(x_1,\ldots,x_r)=
 \sum_{\mathcal P}\frac{(n)_{|\mathcal P|}}{n^r}
 g(x_{\beta_{\mathcal P}(1)},\ldots,x_{\beta_{\mathcal P}(r)}).
\]
The coefficients sum to one and
\(W=U_r(g^*)\), the ordered \(U\) statistic with kernel \(g^*\).
Thus \(a\leq g\leq b\), \(a<b\), on all tuples implies, for \(t>0\),
\(\Pr\{W-EW\geq t\}\leq e^{-2\lfloor n/r\rfloor t^2/(b-a)^2}\).
For \(r=2\) one may use
\[
 g^*(x_1,x_2)=\frac{n-1}{n}g(x_1,x_2)+\frac1n g(x_1,x_1).
\]
For \(r=3\) one valid choice is
\[
\begin{split}
g^*={}&\frac{(n-1)(n-2)}{n^2}g(x_1,x_2,x_3)\\
&+\frac{n-1}{n^2}\{g(x_1,x_1,x_2)+g(x_1,x_2,x_1)+g(x_2,x_1,x_1)\}
+\frac1{n^2}g(x_1,x_1,x_1).
\end{split}
\]
These retain the source's equations (5.16)--(5.17).

A kernel \(g^*\) with \(W=U_r(g^*)\) is not unique. If \(f\) is a kernel
of \(j\leq r\) arguments and \(\iota\) is any injective map of
\(\{1,\ldots,j\}\) into \(\{1,\ldots,r\}\), the kernel
\(f(x_{\iota(1)},\ldots,x_{\iota(j)})\) of \(r\) arguments has the same
ordered average as \(f\). Each partition term may therefore be placed in
any argument positions, or replaced by a convex combination of such
placements. Hoeffding's instance replaces \(g(x_1,x_1)\) in (5.16) by
\(\frac12g(x_1,x_1)+\frac12g(x_2,x_2)\)
\cite[p.~27]{Hoeffding1963BoundedSums}: the kernel
\[
 \tilde g^*(x_1,x_2)=\frac{n-1}{n}g(x_1,x_2)
 +\frac1{2n}\{g(x_1,x_1)+g(x_2,x_2)\}
\]
also satisfies \(W=U_2(\tilde g^*)\) and \(a\leq\tilde g^*\leq b\).
If \(g\) is symmetric then \(\tilde g^*\) is symmetric, which (5.16) need
not be. The choice matters for the bound, since the range of the chosen
kernel can be narrower than that of \(g\); the next theorem computes this
for the source's example.
\end{hbtheorem}

\begin{hbtheorem}{thm:p15-cramer-von-mises-kernel}{proved=annex:p15.cramer-von-mises-kernel}{The Cram\'er--von Mises statistic as a repeated-index example}
For an integer \(n\geq1\), a continuous cumulative distribution function
\(G\) on \(\mathbb R\)
and empirical distribution \(F_n(z)=n^{-1}\sum_i{\bf1}_{\{X_i\leq z\}}\),
define
\[
 \omega^2=\int_{\mathbb R}(F_n(z)-G(z))^2\,dG(z).
\]
This \(\omega^2\), the source's (5.13), is the
\emph{Cram\'er--von Mises goodness-of-fit statistic} for the
hypothesized distribution function \(G\)
\cite[p.~26]{Hoeffding1963BoundedSums}.
\index{Cramer--von Mises statistic@Cram\'er--von Mises statistic}
For every observed sample,
\(\omega^2=n^{-2}\sum_{i,j=1}^n g(X_i,X_j)\), where
\[
 g(x,y)=\frac13+\frac12G(x)^2+\frac12G(y)^2-\max\{G(x),G(y)\}.
\]
For \(n\geq2\) this is the preceding repeated-index statistic of order two.
Its kernel range lies in \([-1/6,1/3]\). For independent \(X_i\),
\(n\geq2\), and \(t>0\),
\(\Pr\{\omega^2-E\omega^2\geq t\}\leq e^{-8\lfloor n/2\rfloor t^2}\).
This bound does not require \(X_i\) to have distribution \(G\).
For two observations with transformed values \(1/4,3/4\),
the statistic is \(1/48\).

The kernel (5.16) built from this \(g\),
\(g^*(x_1,x_2)=\frac{n-1}ng(x_1,x_2)+\frac1ng(x_1,x_1)\), has a
narrower range. For \(n\geq2\) and all real \(x_1,x_2\),
\[
 \frac{2-n}{6(n+1)}\leq g^*(x_1,x_2)\leq\frac13,
\]
and these two numbers are the infimum and supremum of \(g^*\), as
\(-1/6\) and \(1/3\) are those of \(g\). The width is
\(n/\{2(n+1)\}\) instead of \(1/2\), so
\(\Pr\{\omega^2-E\omega^2\geq t\}
 \leq e^{-8\lfloor n/2\rfloor(1+1/n)^2t^2}\).
The width ratio \(n/(n+1)\) tends to one, so the gain is of relative
order \(1/n\) in the exponent; this is the source's remark that the
range of \(g^*\) is smaller than that of \(g\) while the difference is
negligible for large \(n\) \cite[p.~27]{Hoeffding1963BoundedSums}.
At \(n=2\) the range of \(g^*\) is \([0,1/3]\), of width \(1/3\).
\end{hbtheorem}

\begin{hbdefinition}{def:p15-finite-range-dependence}{Dependence between separated blocks}
A finite sequence of real random variables \(Y_1,\ldots,Y_n\)
is \((r-1)\)-dependent if
the vectors \((Y_1,\ldots,Y_i)\) and \((Y_j,\ldots,Y_n)\) are independent
whenever \(j-i\geq r\). Take \(1\leq r\leq n\).
If a stated dependence range has \(r>n\), replace it by \(r=n\);
both conditions then impose no restriction.
Let
\[
 I_i=\{i,i+r,i+2r,\ldots\}\cap\{1,\ldots,n\},\quad
 n_i=|I_i|=\left\lfloor\frac{n-i+r}{r}\right\rfloor,\quad
 S_i=\sum_{j\in I_i}Y_j,\quad k=\lfloor n/r\rfloor .
\]
Hoeffding's example of such a sequence is
\(S=X_1X_r+X_2X_{r+1}+\cdots+X_nX_{r+n-1}\) with independent \(X_i\),
that is, \(Y_i=X_iX_{i+r-1}\) \cite[p.~27]{Hoeffding1963BoundedSums};
its dependence range is verified in
\cref{thm:p15-dependent-concentration-examples}.
\end{hbdefinition}

\begin{hbtheorem}{thm:p15-finite-range-concentration}{proved=annex:p15.finite-range-concentration}{Residue-class decomposition with unequal block sizes}
Assume the preceding dependence condition and integrable \(Y_i\).
Within each \(I_i\) the variables are jointly independent, and, for
\(S=\sum_jY_j\),
\[
 \frac{S-ES}{n}
 =\sum_{i=1}^r\frac{n_i}{n}\frac{S_i-ES_i}{n_i}.
\]
For every \(h>0\) and \(t>0\),
\[
 \Pr\{S-ES\geq nt\}
 \leq\sum_{i=1}^r\frac{n_i}{n}e^{-ht}
       E\exp\{h(S_i-ES_i)/n_i\}.
\]
If \(a\leq Y_i\leq b\), \(a<b\), the bound is
\(e^{-2kt^2/(b-a)^2}\). If in addition the \(Y_i\) have a common
distribution, the common-range bounds for \(0\leq Y_i\leq1\), or
the variance bounds for \(Y_i-EY_i\leq b_0\) with finite variance,
transfer with effective count \(k\).
This remains true when the positive sizes \(n_i\) are unequal.
\end{hbtheorem}

\section{Sampling without replacement}
\label{sec:p15-finite-population-concentration}

\begin{hbdefinition}{def:p15-finite-population-sampling}{The two finite-population sampling laws}
Fix an integer \(N\geq1\) and real values \(c_1,\ldots,c_N\), allowing
repetitions of values.
A sample \(X_1,\ldots,X_n\), \(1\leq n\leq N\), without replacement
uses a uniformly chosen ordered tuple of distinct population indices.
A sample \(Y_1,\ldots,Y_n\) with replacement uses independent uniform
indices from \(\{1,\ldots,N\}\). Write
\[
 \mu=\frac1N\sum_jc_j,\qquad
 \sigma_C^2=\frac1N\sum_j(c_j-\mu)^2.
\]
Repetition of a population value is different from repetition of its index.
\end{hbdefinition}

\begin{hbtheorem}{thm:p15-without-replacement-convex}{proved=annex:p15.without-replacement-convex}{Hoeffding's finite-population convex comparison}
For the two sampling laws and every real continuous convex \(f\),
\[
 E f\!\left(\sum_{i=1}^nX_i\right)
 \leq E f\!\left(\sum_{i=1}^nY_i\right).
\]
Consequently the common-range bounds of
\cref{thm:p15-conc-common-range}, the common-interval sum specialization of
\cref{thm:p15-conc-unequal-ranges}, and the variance bounds of
\cref{thm:p15-variance-tail-chain} for an independent sample from this
population also bound the corresponding unweighted without-replacement
sample mean.
Use the population interval, mean \(\mu\), and variance \(\sigma_C^2\)
in those bounds. In the variance theorem use centered population values
\(c_j-\mu\) and an upper bound for those values.

Both sample means have expectation \(\mu\). For \(N>1\),
\[
 \operatorname{Var}(\overline X)
 =\frac{N-n}{N-1}\frac{\sigma_C^2}{n},\qquad
 \operatorname{Var}(\overline Y)=\frac{\sigma_C^2}{n}.
\]
For \(N=1\) both variances are zero. The inequality
\(\operatorname{Var}(\overline X)\leq\operatorname{Var}(\overline Y)\)
needs neither formula: it is the special case
\(f(s)=(s-n\mu)^2\) of the convex comparison, the source's (6.2)
\cite[p.~28]{Hoeffding1963BoundedSums}.

The proof compares \(f(\sum_iX_i)\) with a symmetrized kernel
\(\overline g\) whose expectation under sampling without replacement
equals \(Ef(\sum_iY_i)\). For a sample of size \(n=2\), \(N\geq2\), it is
\[
 \overline g(x_1,x_2)=\frac{N-1}Nf(x_1+x_2)
 +\frac1{2N}f(2x_1)+\frac1{2N}f(2x_2),
\]
Hoeffding's displayed instance \cite[p.~29]{Hoeffding1963BoundedSums};
by Jensen it is at least \(f(x_1+x_2)\), and
\(E\overline g(X_1,X_2)=Ef(Y_1+Y_2)\).
The proved exponential comparison
does not replace its parameter \(\sigma_C^2\) by the smaller
finite-population-corrected variance. At \(n=N\) the first sample mean is
deterministic even though the transferred upper bound can be positive.
\end{hbtheorem}

\begin{hbtheorem}{thm:p15-dependent-concentration-examples}{proved=annex:p15.dependent-concentration-examples}{Worked blocks, finite samples and a pairwise-independence failure}
The following calculations illustrate the preceding results. Tail bounds
in these examples use \(t>0\).
\begin{enumerate}
\item Let twenty independent variables take values \(1,-1/4\) with
probabilities \(1/5,4/5\). Their means are zero and variances \(1/4\).
At \(t=1/4\), the exact upper tail and four bounds are
\[
\begin{split}
 \Pr\{\overline X\geq1/4\}
 &=5^{-20}\sum_{j=8}^{20}\binom{20}{j}4^{20-j},\\
 B_1&=2^{16}/3^{12},\qquad B_2=e^5/2^{10},\\
 B_{\rm Bernstein}&=e^{-15/8},\qquad
 B_{\rm Prohorov}=e^{-(5/2)\operatorname{arsinh}(1/2)} .
\end{split}
\]
The exact tail is approximately \(0.03214\); the bounds are approximately
\(0.12332,0.14493,0.15335,0.30028\), respectively.
\item Let \(\Pr\{Z=2^j\}=7\,8^{-j}\), \(j=1,2,\ldots\), and
\(X=7/3-Z\). Then \(EX=0\), \(EX^2=14/9\), and \(X\leq1/3\).
This law meets the one-variable upper-tail premises and has no lower
bound. For every \(h>0\), \(Ee^{-hX}=\infty\); the available upper-tail
moment estimate cannot be reused as a lower-tail estimate.
\item For independent identically distributed observations
\(Z_i=(A_i,B_i)\), the order-two kernel
\(g(Z_i,Z_j)=\operatorname{sign}((A_i-A_j)(B_i-B_j))\), with
\(\operatorname{sign}(0)=0\), has range \([-1,1]\).
Its ordered average is the unadjusted Kendall pair statistic and has upper-tail bound
\(e^{-\lfloor n/2\rfloor t^2/2}\). For the three observations
\((1,1),(2,3),(3,2)\), its value is \(1/3\).
\item For independent samples, the proportion of ordered pairs satisfying
\(Y_j<X_i\) uses \(g(x,y)={\bf1}_{\{y<x\}}\), with \(r=s=1\).
This is the normalized Wilcoxon--Mann--Whitney pair count.
Its upper-tail bound is \(e^{-2\min(m,n)t^2}\).
For \(X=(1,3)\), \(Y=(0,2,4)\), the observed proportion is \(3/6=1/2\).
\item If \(Z_1,\ldots,Z_{n+r-1}\) are independent and
\(Y_i=Z_i\cdots Z_{i+r-1}\), then \(Y_i\) is \((r-1)\)-dependent.
For independent Bernoulli\((p)\) inputs, \(0\leq Y_i\leq1\),
\(EY_i=p^r\), and its empirical mean has the common-range bounds with
effective count \(\lfloor n/r\rfloor\).
For \(n=5,r=2\), the two residue sets have sizes \(3,2\), with weights
\(3/5,2/5\), and the range-only bound is \(e^{-4t^2}\).
\item In the source's example, \(X_1,\ldots,X_{n+r-1}\) are independent,
\(Y_i=X_iX_{i+r-1}\), and \(S=X_1X_r+X_2X_{r+1}+\cdots+X_nX_{r+n-1}\).
The sequence \(Y_1,\ldots,Y_n\) is \((r-1)\)-dependent. For \(r=2\) it is
the preceding moving product; for \(r\geq3\) it differs, since then
\(Y_i\) and \(Y_{i+1}\) share no factor. The range \(r-1\) cannot be
lowered in general: for \(2\leq r\leq n\) and independent
Bernoulli\((p)\) inputs with \(0<p<1\),
\(\Pr\{Y_1=1,Y_r=1\}=p^3\neq p^4=\Pr\{Y_1=1\}\Pr\{Y_r=1\}\), so the
sequence is not \((r-2)\)-dependent. For these inputs and \(r\geq2\),
each \(Y_i\) is Bernoulli\((p^2)\), and
\(\Pr\{S-np^2\geq nt\}\leq e^{-2\lfloor n/r\rfloor t^2}\).
\item From population \((0,0,1,1)\), sample two indices.
The without-replacement sum has probabilities \((1/6,4/6,1/6)\)
at \(0,1,2\); the with-replacement sum has probabilities
\((1/4,1/2,1/4)\). Their mean variances are \(1/12\) and \(1/8\).
The with-replacement exponential moment minus the without-replacement
one is \((e^h-1)^2/12\). At displacement \(t=1/2\), the exact
without-replacement upper tail is \(1/6\); the transferred endpoint
bound is \(1/4\).
\item If \(A,B\) are independent fair random signs, the three variables
\(A,B,AB\) are pairwise independent, centered, upper bounded by one,
and each has variance one. Their mean equals one with probability
\(1/4\), exceeding the independent variance-theorem endpoint bound
\(2^{-3}=1/8\). Pairwise independence does not suffice.
\end{enumerate}
\end{hbtheorem}

The derived bounds are finite-sample probability statements under the
specified joint laws. They do not infer a sampling law from a table of
observations. In particular overlapping kernel terms and adjacent blocks
cannot be counted as mutually independent observations.
